\documentclass[10pt,a4paper]{article}

\usepackage{fontspec}
\usepackage{polyglossia}
\setmainlanguage{russian}
\setmainfont{Arial}

\usepackage{amsmath,amssymb}
\usepackage[margin=1.35cm]{geometry}
\usepackage{enumitem}
\usepackage[hidelinks]{hyperref}

\newcommand{\X}{\mathcal X}
\newcommand{\A}{\mathcal A}
\newcommand{\B}{\mathcal B}
\newcommand{\C}{\mathbb C}
\newcommand{\R}{\mathbb R}
\newcommand{\Lin}{\operatorname{Lin}}
\newcommand{\Ker}{\operatorname{Ker}}
\newcommand{\Ree}{\operatorname{Re}}
\newcommand{\wstar}{w^\ast}

\setlength{\parindent}{0pt}
\setlength{\parskip}{0.25em}
\setlist[enumerate]{leftmargin=1.35em,itemsep=0.15em,topsep=0.1em,parsep=0pt,partopsep=0pt}
\setlist[itemize]{leftmargin=1.25em,itemsep=0.1em,topsep=0.1em,parsep=0pt,partopsep=0pt}

\begin{document}
\small
\sloppy

\begin{center}
    {\large\bfseries Функциональный анализ: пункты 1--29}
\end{center}

\section*{1. Слабая-* топология. Слабо-* непрерывные функционалы}
\textbf{Формулировка.} Для $x\in\X$ рассмотрим $\Phi_x\in(\X^\ast)^\ast$, $\Phi_x(f)=f(x)$. Слабая-* топология на $\X^\ast$ --- слабая топология, порождённая семейством $\Phi(\X)$; то есть это слабейшая топология, в которой все $f\mapsto f(x)$ непрерывны. Сходимость $f_n\overset{\wstar}{\to} f$ равносильна поточечной: $f_n(x)\to f(x)$ для всех $x\in\X$. По лемме о сопряжённом пространстве все слабo-* непрерывные линейные функционалы на $\X^\ast$ имеют вид $F(f)=\sum_{k=1}^n\alpha_k f(x_k)=f(x)$, где $x=\sum\alpha_kx_k$.

\textbf{Доказательство, идея.}
\begin{enumerate}
    \item Конечные комбинации $\Phi_x$ непрерывны по определению топологии.
    \item Если $F$ непрерывен в нуле, то некоторая базовая окрестность по $x_1,\dots,x_n$ лежит в $\{|F|<1\}$.
    \item Отсюда $\bigcap_k\Ker \Phi_{x_k}\subset\Ker F$: если $g$ зануляет все $x_k$, то $tg$ лежит в этой окрестности при всех $t>0$, значит $F(g)=0$.
    \item По лемме о ядрах $F\in\Lin\{\Phi_{x_1},\dots,\Phi_{x_n}\}$.
\end{enumerate}

\section*{2. Банах--Алаоглу}
\textbf{Формулировка.} Пусть $(\X,\tau)$ --- топологическое векторное пространство, $0\in V\in\tau$. Поляра окрестности $V$ есть $\Gamma(V)=\{f\in(\X,\tau)^\ast:\ |f(x)|\le1\ \forall x\in V\}$, где $(\X,\tau)^\ast$ --- топологическое сопряжённое. Тогда $\Gamma(V)$ компактна в слабой-* топологии на $(\X,\tau)^\ast$.

\textbf{Доказательство, идея.}
\begin{enumerate}
    \item Смотрим на $(\X,\tau)^\ast\subset\C^\X$ с топологией Тихонова; её сужение --- слабая-* топология, поэтому достаточно получить компактность в топологии Тихонова и затем перейти к подпространству.
    \item Для каждого $x$ из непрерывности умножения на скаляр в ТВП выбираем $\delta_x>0$ так, что $\delta_xx\in V$. Тогда для $f\in\Gamma(V)$ имеем $|f(x)|\le1/\delta_x$, значит $\Gamma(V)$ лежит в тихоновском компакте $\prod_x\{|\lambda|\le1/\delta_x\}$.
    \item Теперь осталось показать замкнутость поляры: пусть $h\in[\Gamma(V)]_{\tau_T}$. Надо показать, что $h$ линеен, непрерывен и $|h|\le1$ на $V$.
    \item Линейность: для $x,y,\alpha,\beta$ и $\varepsilon>0$ берём $f\in\Gamma(V)$ близкий к $h$ в координатах $x,y,\alpha x+\beta y$. Тогда $|h(\alpha x+\beta y)-\alpha h(x)-\beta h(y)|\le\varepsilon(1+|\alpha|+|\beta|)$, значит равенство точное.
    \item Оценка на $V$: для $x\in V$ берём $f\in\Gamma(V)$ с $|h(x)-f(x)|<\varepsilon$; тогда $|h(x)|\le1+\varepsilon$, значит $|h(x)|\le1$.
    \item Непрерывность: из $|h|\le1$ на $V$ следует $|h|\le\varepsilon$ на $\varepsilon V$, то есть $h$ непрерывен в нуле. Значит $h\in(\X,\tau)^\ast$ и $h\in\Gamma(V)$; поляра замкнута в компакте.
\end{enumerate}

\section*{3/9. Метризуемость слабых топологий на шаре}
\textbf{Формулировка.} Два одинаковых по идее случая. Если $\X$ сепарабельно, то слабая-* топология на шаре $\B_R^\ast[0]\subset\X^\ast$ метризуема; отсюда этот шар слабo-* секвенциально компактен. Если $\X^\ast$ сепарабельно, то слабая топология на шаре $\B_R[0]\subset\X$ метризуема.

\textbf{Метрика.} В слабo-* случае берём плотное $\{x_n\}$ на единичной сфере $\X$ и кладём $\rho_\ast(f,g)=\sum_{n\ge1}2^{-n}|(f-g)(x_n)|$. В слабом случае берём плотное $\{f_n\}$ на единичной сфере $\X^\ast$ и кладём $\rho(x,y)=\sum_{n\ge1}2^{-n}|f_n(x-y)|$.

\textbf{Доказательство, идея.}
\begin{enumerate}
    \item Это метрика: нулевое расстояние даёт совпадение на плотном множестве координат; дальше по непрерывности и разделимости точек функционалами.
    \item Включение ``метрика $\Rightarrow$ слабая топология'': произвольную координату приближаем плотной координатой; ошибка оценивается ограниченностью шара, например в слабo-* случае $\|f-g\|\le2R$.
    \item Обратное включение: чтобы попасть в метрический шар радиуса $r$, контролируем конечное число первых координат, а хвост ряда делаем малым выбором $N$.
    \item Различие только в том, что в слабo-* случае координаты --- точки $x_n\in\X$, а в слабом --- функционалы $f_n\in\X^\ast$.
    \item Секвенциальная компактность слабo-* шара: по Банаху--Алаоглу он слабo-* компактен, а на нём топология метрическая.
\end{enumerate}

\section*{4/8. Неметризуемость слабых топологий}
\textbf{Формулировка.} Если $\X$ бесконечномерно нормированно, то слабая топология на $\X$ не метризуема. Если $\X$ бесконечномерно банахово, то слабая-* топология на $\X^\ast$ не метризуема.

\textbf{Доказательство, идея.}
\begin{enumerate}
    \item \textit{Слабая топология на $\X$.} Допустим, метрика есть; у нуля есть счётная база $O_{1/n}(0)$. Каждый $O_{1/n}$ содержит $\bigcap_{k=1}^{m_n}\Ker f_{k,n}$ для некоторых $f_{k,n}\in\X^\ast$.
    \item Для произвольного $f\in\X^\ast$ окрестность $\{|f(x)|<1\}$ содержит некоторое $O_{1/n}$. Тогда $\bigcap_k\Ker f_{k,n}\subset\Ker f$: если $x$ лежит в пересечении ядер, то и $tx$ лежит там при всех $t>0$, значит $|f(tx)|<1$ при всех $t$, откуда $f(x)=0$.
    \item По лемме о ядрах $f\in L_n=\Lin\{f_{k,n}\}$. Получаем счётное объединение $\X^\ast=\bigcup_nL_n$, где $L_n$ конечномерны, а значит замкнуты. По Бэру в банаховом $\X^\ast$ одно $L_n$ имеет внутренность, значит $\X^\ast$ конечномерно, а значит и $\X$; противоречие.
    \item \textit{Слабая-* топология на $\X^\ast$.} То же рассуждение, но элементы пространства теперь функционалы $g\in\X^\ast$, а координаты задаются точками $x\in\X$: базовая окрестность содержит $\bigcap_k\Ker\Phi_{x_{k,n}}$.
    \item Для произвольного $x\in\X$ окрестность $\{|g(x)|<1\}$ даёт $\bigcap_k\Ker\Phi_{x_{k,n}}\subset\Ker\Phi_x$ тем же трюком с $tg$. По лемме о ядрах $x\in\Lin\{x_{k,n}\}$, значит $\X=\bigcup_n\Lin\{x_{k,n}\}$.
    \item Здесь Бэр применяем уже к банахову $\X$: одно конечномерное подпространство имеет внутренность, значит равно всему $\X$. Противоречие бесконечномерности.
\end{enumerate}

\section*{5. Первая теорема об отделимости}
\textbf{Формулировка.} Пусть $(\X,\tau)$ --- топологическое векторное пространство, $\A,\B$ --- непустые выпуклые множества, $\A\cap\B=\varnothing$, причём $\A$ открыто. Тогда существуют $f\in(\X,\tau)^\ast$ и $r\in\R$: для всех $a\in\A$, $b\in\B$ выполнено $\Ree f(a)<r\le\Ree f(b)$.

\textbf{Доказательство, идея.}
\begin{enumerate}
    \item Рассматриваем $C=\A-\B$: оно открыто, выпукло и не содержит $0$. Берём $c_0=a_0-b_0\in C$ и $G=C-c_0$; тогда $G$ --- выпуклая окрестность нуля, а $-c_0\notin G$.
    \item Для $G$ берём функционал Минковского $\mu_G(x)=\inf\{t>0:x/t\in G\}$; он положительно однороден и полуаддитивен.
    \item На вещественной прямой $L=\Lin_\R\{c_0\}$ задаём $\varphi(\lambda c_0)=-\lambda$. Так как $\mu_G(-c_0)\ge1$, имеем $\varphi\le\mu_G$ на $L$.
    \item По Хану--Банаху продолжаем до вещественно-линейного $\psi$ на $\X$ с $\psi\le\mu_G$.
    \item Для $a\in\A$, $b\in\B$: $a-b-c_0\in G$, значит $\psi(a)-\psi(b)+1=\psi(a-b-c_0)\le1$, то есть $\psi(a)\le\psi(b)$.
    \item Непрерывность $\psi$: на $G\cap(-G)$ получаем $|\psi|\le1$, затем масштабируем окрестность нуля.
    \item Комплексифицируем: $f(x)=\psi(x)-i\psi(ix)$, тогда $\Ree f=\psi$.
    \item Строгость: берём $r=\sup_{\A}\psi$. Открытость $\A$ позволяет сдвинуть $a$ немного в направлении $-c_0$, поэтому $\psi(a)+\delta=\psi(a-\delta c_0)\le r$, значит $\psi(a)<r$. А из $\psi(a)\le\psi(b)$ следует $r\le\psi(b)$.
\end{enumerate}

\section*{6. Вторая теорема об отделимости}
\textbf{Формулировка.} В локально выпуклом ТВП $(\X,\tau)$, если $M\subset\X$ непусто, выпукло и замкнуто, а $x_0\notin M$, то существуют $f\in(\X,\tau)^\ast$ и $r\in\R$: $\Ree f(x_0)>r\ge\Ree f(x)$ для всех $x\in M$. Локальная выпуклость означает: у нуля есть база, состоящая из выпуклых окрестностей.

\textbf{Доказательство, идея.}
\begin{enumerate}
    \item Так как $M$ замкнуто и $x_0\notin M$, точка $x_0$ не является точкой прикосновения $M$, значит есть окрестность $U(x_0)$, не пересекающая $M$.
    \item Сдвигом получаем окрестность нуля $U(x_0)-x_0$; по локальной выпуклости берём выпуклую $V(0)\subset U(x_0)-x_0$.
    \item Тогда $A=x_0+V(0)$ --- открытая выпуклая окрестность $x_0$; сдвиги в ТВП сохраняют открытость, поэтому такая выпуклая окрестность есть и у $x_0$, и она не пересекает $M$.
    \item Применяем первую теорему об отделимости к $A$ и $M$.
\end{enumerate}

\section*{7. Слабая топология, сравнение с нормированной. Мазур}
\textbf{Формулировка.} В нормированном $\X$ слабая топология $\sigma(\X,\X^\ast)$ задаётся всеми $f\in\X^\ast$; $x_n\rightharpoonup x$ iff $f(x_n)\to f(x)$ для всех $f$. Она слабее нормированной, а при $\dim\X=\infty$ строго слабее. Теорема Мазура: выпуклое нормозамкнутое $M\subset\X$ слабо замкнуто.

\textbf{Доказательство, идея.}
\begin{enumerate}
    \item Слабая топология слабее нормированной, потому что все $f\in\X^\ast$ нормово непрерывны.
    \item При бесконечной размерности любая слабая окрестность содержит $x+\bigcap_{k=1}^n\Ker f_k$, то есть сдвиг бесконечномерного подпространства; поэтому она не может лежать в малом нормовом шаре.
    \item Для Мазура берём $x\notin M$. Вторая теорема об отделимости даёт $f\in\X^\ast$: $\Ree f(x)>r\ge\Ree f(y)$ на $M$.
    \item Слабая окрестность $x$ по функционалу $f$ не пересекает $M$, значит дополнение $M$ слабо открыто.
\end{enumerate}

\section*{10. Сепарабельность пространства из сепарабельности сопряжённого}
\textbf{Формулировка.} Если нормированное $\X^\ast$ сепарабельно, то $\X$ сепарабельно.

\textbf{Доказательство, идея.}
\begin{enumerate}
    \item Берём плотное $\{f_n\}\subset\X^\ast$; для каждого ненулевого $f_n$ выбираем $\|x_n\|=1$, почти достигающий норму: $\|f_n\|\le2|f_n(x_n)|$.
    \item Пусть $L=\overline{\Lin\{x_n\}}$. Замыкание нужно, чтобы получить замкнутое подпространство; при этом рационально-комплексные комбинации $x_n$ счётны и плотны в $L$, значит $L$ сепарабельно.
    \item Если $L\ne\X$, берём $z\notin L$ и задаём на $L\oplus\Lin\{z\}$ функционал $\varphi|_L=0$, $\varphi(z)=1$; по Хану--Банаху продолжаем его до ненулевого $f\in\X^\ast$, зануляющего $L$.
    \item Приближаем $f$ некоторым $f_n$: $\|f-f_n\|<\varepsilon$. Так как $x_n\in L$, имеем $f(x_n)=0$, поэтому $\|f_n\|\le2|f_n(x_n)|=2|(f_n-f)(x_n)|\le2\varepsilon$.
    \item Тогда $\|f\|\le\|f-f_n\|+\|f_n\|\le3\varepsilon$ при произвольном $\varepsilon$, значит $f=0$, противоречие с $\varphi(z)=1$.
\end{enumerate}

\section*{11. Эберлейн--Шмульян}
\textbf{Формулировка.} Слабый компакт $K$ в нормированном пространстве слабo секвенциально компактен.

\textbf{Доказательство, идея.}
\begin{enumerate}
    \item Сначала $K$ нормово ограничен: для максимизирующей последовательности норм все $f(x_n)$ ограничены, и Банах--Штейнгауз даёт ограниченность $\|x_n\|$.
    \item Для последовательности $\{x_n\}\subset K$ берём сепарабельное $L=\overline{\Lin\{x_n\}}$; $K\cap L$ остаётся слабым компактом.
    \item Так как $L$ сепарабельно, шар в $L^\ast$ слабo-* метризуем и компактен, значит в нём есть счётное плотное $\{f_k\}$.
    \item Диагональным процессом выделяем подпоследовательность, на которой все $f_k(x_{n_m})$ сходятся.
    \item Компактность даёт предельную точку $x_0\in K\cap L$; плотность $\{f_k\}$ в шаре $L^\ast$ переносит сходимость на все функционалы, значит $x_{n_m}\rightharpoonup x_0$.
\end{enumerate}

\section*{12. Слабая компактность шара в рефлексивном пространстве. Проекция}
\textbf{Формулировка.} Если $\X$ рефлексивно, то замкнутый единичный шар слабo компактен. Если $M\subset\X$ непусто, выпукло и сильно замкнуто, то для любого $z\in\X$ существует элемент наилучшего приближения к $z$ в $M$.

\textbf{Доказательство, идея.}
\begin{enumerate}
    \item В рефлексивном случае $\X=J_\X\X\subset\X^{\ast\ast}$; слабая топология на $\X$ соответствует слабой-* топологии на $J_\X\X$.
    \item По Банаху--Алаоглу шар в $\X^{\ast\ast}$ слабo-* компактен, значит шар в $\X$ слабo компактен.
    \item Для проекции минимизируем $J(x)=\|x-z\|$ на $M$. Подуровень по точке $y_0\in M$ ограничен, выпукл и замкнут.
    \item На этом подуровне применяем слабую компактность и Мазура; выпуклость и сильная непрерывность $J$ дают достижение инфимума.
\end{enumerate}

\section*{13. Слабая компактность шара как критерий рефлексивности}
\textbf{Формулировка.} Банахово $\X$ рефлексивно тогда и только тогда, когда его замкнутый единичный шар слабo компактен.

\textbf{Доказательство, идея.}
\begin{enumerate}
    \item Направление ``рефлексивность $\Rightarrow$ компактность'' --- предыдущий пункт.
    \item Обратно, если $\B_1[0]\subset\X$ слабo компактна, то $J_\X\B_1[0]$ слабo-* компактен в $\X^{\ast\ast}$, значит слабo-* замкнут.
    \item Если бы существовал $\Phi\in \B_1^{\ast\ast}\setminus J_\X\B_1$, отделяем $\Phi$ от $J_\X\B_1$ слабo-* непрерывным функционалом, то есть элементом $f\in\X^\ast$.
    \item Получается невозможная оценка: $\Re\Phi(f)>r\ge\sup_{\|x\|\le1}\Re f(x)=\|f\|$ при $\|\Phi\|\le1$.
    \item Значит $J_\X\B_1=\B_1^{\ast\ast}$; масштабированием $J_\X\X=\X^{\ast\ast}$.
\end{enumerate}

\section*{14. Минимум выпуклого функционала}
\textbf{Формулировка.} В рефлексивном $\X$ сильно непрерывный выпуклый вещественный функционал $J$ достигает минимума на выпуклом сильно замкнутом ограниченном $M$.

\textbf{Доказательство, идея.}
\begin{enumerate}
    \item Берём минимизирующую последовательность $x_n\in M$.
    \item $M$ слабo компактен и, по Эберлейну--Шмульяну, слабo секвенциально компактен; выделяем $x_{n_k}\rightharpoonup x\in M$.
    \item По следствию Мазура есть выпуклые комбинации хвоста $y_m\in\operatorname{conv}\{x_{n_k}:k\ge K\}$, сходящиеся к $x$ по норме.
    \item Выпуклость даёт $J(y_m)\le\gamma$ для любого $\gamma>\inf_MJ$ при достаточно далёком хвосте; сильная непрерывность даёт $J(x)\le\gamma$.
    \item Пускаем $\gamma\downarrow\inf_MJ$.
\end{enumerate}

\section*{15. Рисс--Фреше. Рефлексивность гильбертова пространства}
\textbf{Формулировка.} Для гильбертова $\mathcal H$ всякий $f\in\mathcal H^\ast$ единственно представим как $f(x)=\langle x,z_f\rangle$, причём $\|f\|=\|z_f\|$. Всякое гильбертово пространство рефлексивно.

\textbf{Доказательство, идея.}
\begin{enumerate}
    \item Для $f\ne0$ берём разложение $\mathcal H=\Ker f\oplus(\Ker f)^\perp$ и ненулевой $z\in(\Ker f)^\perp$.
    \item Любой $x$ раскладывается как $x=f(x)z/f(z)+y$, $y\in\Ker f$; отсюда после нормировки $f(x)=\langle x,z_f\rangle$.
    \item Единственность: если $\langle x,z_1-z_2\rangle=0$ для всех $x$, то $z_1=z_2$; норма --- из КБШ и подстановки $x=z_f/\|z_f\|$.
    \item Для рефлексивности берём $F\in\mathcal H^{\ast\ast}$, переводим его через изометрию Рисса в функционал на $\mathcal H$, снова применяем Рисс--Фреше и получаем $F(f)=f(y)$.
\end{enumerate}

\section*{16. Аннуляторы и теорема Голдстайна}
\textbf{Формулировка.} Для $L\subset\X$, $N\subset\X^\ast$: $L^\perp=\{f:f|_L=0\}$, ${}^\perp N=\{x:f(x)=0\ \forall f\in N\}$. Верно ${}^\perp(L^\perp)=\overline L^{\|\cdot\|}=\overline L^w$ и $({}^\perp N)^\perp=\overline N^{w^\ast}\supset\overline N^{\|\cdot\|}$. Теорема Голдстайна: $J_\X(\X)$ слабo-* плотно в $\X^{\ast\ast}$.

\textbf{Доказательство, идея.}
\begin{enumerate}
    \item Первое равенство: включение $\overline L\subset{}^\perp(L^\perp)$ очевидно; если $x\notin\overline L$, отделяем $x$ от $\overline L$ функционалом, который зануляет $L$.
    \item Второе: работаем в слабой-* локально выпуклой топологии на $\X^\ast$; если $f\notin\overline N^{w^\ast}$, отделяем его слабo-* непрерывным функционалом, а такие имеют вид $J_x$.
    \item Тогда $J_x$ зануляет $N$, то есть $x\in{}^\perp N$, но не зануляет $f$ --- противоречие принадлежности $f\in({}^\perp N)^\perp$.
    \item Голдстайн: $\overline{J_\X\X}^{w^\ast}=({}^\perp J_\X\X)^\perp$, а ${}^\perp J_\X\X=\{0\}\subset\X^\ast$, значит замыкание равно $\X^{\ast\ast}$.
\end{enumerate}

\section*{17. Рефлексивность $\X$ и $\X^\ast$}
\textbf{Формулировка.} Для банахова $\X$: $\X$ рефлексивно тогда и только тогда, когда $\X^\ast$ рефлексивно.

\textbf{Доказательство, идея.}
\begin{enumerate}
    \item Если $\X^\ast$ рефлексивно, доказываем $J_\X\X=\X^{\ast\ast}$ через аннулятор: достаточно $J_\X(\X)^\perp=\{0\}$ в $\X^{\ast\ast\ast}$.
    \item Любой $\Phi\in\X^{\ast\ast\ast}$ имеет вид $J_{\X^\ast}f$; если он зануляет все $J_\X x$, то $f(x)=0$ для всех $x$, значит $f=0$.
    \item Если $\X$ рефлексивно, то для $\Phi\in\X^{\ast\ast\ast}$ берём $f=J_\X^\ast\Phi$ и проверяем $J_{\X^\ast}f=\Phi$ на $J_\X\X=\X^{\ast\ast}$.
\end{enumerate}

\section*{18. Сопряжённый оператор и аннуляторные равенства}
\textbf{Формулировка.} Для $A\in\mathcal L(\X,Y)$: $A^\ast g=g\circ A$, $\|A^\ast\|=\|A\|$. Верно $\Ker A={}^\perp\operatorname{Im}A^\ast$, $\Ker A^\ast=(\operatorname{Im}A)^\perp$, поэтому ${}^\perp(\Ker A^\ast)=\overline{\operatorname{Im}A}^{\|\cdot\|}$ и $(\Ker A)^\perp=\overline{\operatorname{Im}A^\ast}^{w^\ast}$. Если $\X,Y$ банаховы и $\operatorname{Im}A$ замкнут, то $(\Ker A)^\perp=\operatorname{Im}A^\ast$.

\textbf{Доказательство, идея.}
\begin{enumerate}
    \item Первые равенства --- прямое раскрытие: $(A^\ast g)(x)=g(Ax)$.
    \item Замыкания получаются применением формул для двойных аннуляторов.
    \item При замкнутом образе для $f\in(\Ker A)^\perp$ задаём $\varphi(Ax)=f(x)$ на $\operatorname{Im}A$.
    \item Корректность идёт из зануления $\Ker A$; непрерывность --- из теоремы об открытом отображении для $A:\X\to\operatorname{Im}A$.
    \item Продлеваем $\varphi$ на $Y$ по Хану--Банаху: $g|_{\operatorname{Im}A}=\varphi$, тогда $f=A^\ast g$.
\end{enumerate}

\section*{19. Замкнутость $\operatorname{Im}A$ из замкнутости $\operatorname{Im}A^\ast$}
\textbf{Формулировка.} Если $\X,Y$ банаховы, $A\in\mathcal L(\X,Y)$ и $\operatorname{Im}A^\ast$ замкнут в $\X^\ast$, то $\operatorname{Im}A$ замкнут и $\operatorname{Im}A={}^\perp\Ker A^\ast$, $\operatorname{Im}A^\ast=(\Ker A)^\perp$.

\textbf{Доказательство, идея.}
\begin{enumerate}
    \item Кладём $Z=\overline{\operatorname{Im}A}$ и рассматриваем $T:\X\to Z$, $Tx=Ax$; тогда $\operatorname{Im}T$ плотно в $Z$.
    \item У $T^\ast$ ядро нулевое, а $\operatorname{Im}T^\ast=\operatorname{Im}A^\ast$ замкнут и полон.
    \item Значит $(T^\ast)^{-1}:\operatorname{Im}T^\ast\to Z^\ast$ ограничен.
    \item Через отделимость показываем, что шар в $Z$ лежит в замыкании $T$-образа шара; затем стандартной итерацией с радиусами $1,1/2,1/4,\dots$ убираем замыкание и получаем шар в самом $\operatorname{Im}T$.
    \item Следовательно $\operatorname{Im}T=Z$, то есть $\operatorname{Im}A$ замкнут; равенства уже следуют из пункта 18.
\end{enumerate}

\section*{20. Ядро и образ $A-\lambda I$ для компактного оператора}
\textbf{Формулировка.} Если $\X$ банахово, $A \in \mathcal{L}(\X)$ компактен (т.е. $A\mathcal{B}_1(0)$ вполне ограничен), $\lambda\in\mathbb{C}\setminus\{0\}$, $A_\lambda=A-\lambda I$, то $\dim\Ker A_\lambda<\infty$ и $\operatorname{Im}A_\lambda$ замкнут.

\textbf{Доказательство, идея.}
\begin{enumerate}

    \item Хотим показать, что в $\Ker A_\lambda$ единичный шар вполне ограничен: если $\{ x_n \} \subset \Ker A_\lambda$, то $x_n=Ax_n/\lambda \in \frac 1 \lambda AB_R(0)$. $A$ компактно, значит $AB_R(0)$ вполне ограничено и в $\{ x_n \}$ есть фундаментальная подпоследовательность.
    \item Образ единичного шара вполне ограничен, значит $\dim \Ker A < \infty$.
    \item Берём замкнутое дополнение: $X=\Ker A_\lambda \oplus N$. $N = \bigcap f_k$, где $f_k - $ координатные функционалы $\Ker A_\lambda \rightarrow \mathbb{C}$, продолженные по теореме Хана-Банаха на  $X$.

    \item Доказываем оценку снизу на $N$: $\exists c>0:\ \|A_\lambda x\|\ge c\|x\|$ для всех $x\in N$. Иначе $\exists x_n\in N$, $\|x_n\|=1$, такие что $A_\lambda x_n\to0$.

    \item По компактности $A$ выделяем $Ax_{n_k}\to y$. Тогда $x_{n_k}=\lambda^{-1}(Ax_{n_k}-A_\lambda x_{n_k})$ сходится к некоторому $x\in N$.

    \item Но $A_\lambda x=0$, значит $x\in\Ker A_\lambda$. Поэтому $x\in N\cap\Ker A_\lambda=\{0\}$, хотя $\|x\|=1$. Противоречие, оценка доказана.

    \item Пусть $y_n=A_\lambda z_n\to y$. Пишем $z_n=u_n+v_n$, где $u_n\in\Ker A_\lambda$, $v_n\in N$. Тогда $y_n=A_\lambda v_n$.

    \item По оценке снизу $v_n$ фундаментальна, значит $v_n\to v\in N$. Следовательно, $y=\lim A_\lambda v_n=A_\lambda v$, то есть $y\in \text{Im} A_\lambda$. Значит $\text{Im} A_\lambda$ замкнут.

\end{enumerate}

\section*{21. Компактность $A$ и $A^\ast$}
\textbf{Формулировка.} Пусть $\mathcal X$ --- нормированное пространство, $\mathcal Y$ --- банахово пространство, $\mathcal A\in\mathcal L(\mathcal X,\mathcal Y)$. Тогда $\mathcal A$ компактен тогда и только тогда, когда $\mathcal A^*\in\mathcal L(\mathcal Y^*,\mathcal X^*)$ компактен.

\textbf{Доказательство, идея.}
\begin{enumerate}
    \item Пусть $\mathcal A$ компактен. Тогда $K=\overline{\mathcal A(\mathcal B_1[0])}$ компактен в $\mathcal Y$, так как $\mathcal A(\mathcal B_1[0])$ вполне ограничено, а $\mathcal Y$ полно.

    \item Берём последовательность $g_n\in \mathcal B_1^{\mathcal Y^*}[0]$ и рассматриваем ограничения $g_n|_K\in C(K)$.

    \item Семейство $\{g_n|_K\}$ равномерно ограничено и равностепенно непрерывно: $|g_n(y)-g_n(z)|\le \|y-z\|$ для всех $y,z\in K$.

    \item По теореме Арцела--Асколи выделяем подпоследовательность $g_{n_k}|_K$, сходящуюся в $C(K)$.

    \item Тогда $\|\mathcal A^*g_{n_k}-\mathcal A^*g_{n_m}\|=\sup_{\|x\|\le1}|g_{n_k}(\mathcal A x)-g_{n_m}(\mathcal A x)|\le \|g_{n_k}|_K-g_{n_m}|_K\|_{C(K)}$, значит $\mathcal A^*g_{n_k}$ фундаментальна в $\mathcal X^*$, то есть $\mathcal A^*$ компактен.

    \item Обратно, пусть $\mathcal A^*$ компактен. По уже доказанному для $\mathcal A^*\in\mathcal L(\mathcal Y^*,\mathcal X^*)$ получаем, что $\mathcal A^{**}\in\mathcal L(\mathcal X^{**},\mathcal Y^{**})$ компактен.

    \item Пусть $\Phi_{\mathcal X}:\mathcal X\to\mathcal X^{**}$ и $\Phi_{\mathcal Y}:\mathcal Y\to\mathcal Y^{**}$ --- канонические изометрии. Тогда для всех $x\in\mathcal X$ верно $\mathcal A^{**}\Phi_{\mathcal X}(x)=\Phi_{\mathcal Y}(\mathcal A x)$.

    \item Поэтому $\Phi_{\mathcal Y}(\mathcal A(\mathcal B_1[0]))=\mathcal A^{**}(\Phi_{\mathcal X}(\mathcal B_1[0]))$ вполне ограничено в $\mathcal Y^{**}$.

    \item Так как $\Phi_X, \Phi_{\mathcal Y}$ --- изометрии, значит сохраняется вполне ограниченность. Значит $\mathcal A$ компактен.
\end{enumerate}

\section*{22. Равенство размерностей ядер в гильбертовом случае}
\textbf{Формулировка.} $H$ --- гильбертово, $A\in\mathcal L(H)$ компактен, $\lambda\in\mathbb C\setminus\{0\}$, $A_\lambda=A-\lambda I$. Тогда $\dim\ker A_\lambda=\dim\ker(A_\lambda)^*$.

\textbf{Доказательство, идея.}
\begin{enumerate}
    \item В силу аннуляторных соотношений, $\ker(A_\lambda)^*=(\text{Im}\,A_\lambda)^\perp, \Ker A_{\lambda} = (\text{Im}(A_{\lambda})^*)^\perp$.

    \item Достаточно доказать $\dim\ker A_\lambda\le \dim(\text{Im}\,A_\lambda)^\perp$, а обратное неравенство затем применить к $(A_\lambda)^*$.

    \item Предположим противное: $\dim\ker A_\lambda>\dim(\text{Im}\,A_\lambda)^\perp$. Тогда существует сюръективный оператор $F:\ker A_\lambda\to(\text{Im}\,A_\lambda)^\perp$ с нетривиальным ядром.

    \item Так как $\ker A_\lambda$ замкнуто, по теореме Рисса о разложении существует ортопроектор $P:H\to\ker A_\lambda$.

    \item Рассмотрим $T=A+FP$. Оператор $FP$ конечномерный, значит компактный, поэтому $T$ компактен.

    \item Возьмём $0\ne x_0\in\ker F\subset\ker A_\lambda$. Тогда $Px_0=x_0$ и $T_\lambda x_0=A_\lambda x_0+FPx_0=0$, то есть $\ker T_\lambda\ne\{0\}$.

    \item По предыдущей теореме для компактного оператора из $\ker T_\lambda\ne\{0\}$ следует, что $\text{Im}\,T_\lambda\ne H$.

    \item Но $\text{Im}\,T_\lambda=T_\lambda(\ker A_\lambda\oplus(\ker A_\lambda)^\perp)=A_\lambda((\ker A_\lambda)^\perp)+F(\ker A_\lambda)=\text{Im}\,A_\lambda\oplus(\text{Im}\,A_\lambda)^\perp=H$. Противоречие.

    \item Значит $\dim\ker A_\lambda\le\dim\ker(A_\lambda)^*$. Применяя то же рассуждение к $(A_\lambda)^*$, получаем обратное неравенство, поэтому размерности равны.
\end{enumerate}

\section*{23. Критерий Фредгольма. Альтернатива Фредгольма}
\textbf{Формулировка.} Пусть $X$ --- банахово пространство, $A\in\mathcal L(X)$ компактен, $\lambda\in\mathbb C\setminus\{0\}$, $A_\lambda=A-\lambda I$. \newline Тогда $\forall v\in X\quad (\exists u\in X:\ A_\lambda u=v)\Longleftrightarrow(\forall f\in\ker A_\lambda^*\quad f(v)=0)$. \newline Альтернатива Фредгольма: либо $\forall v\in X\ \exists!u\in X:\ A_\lambda u=v$, либо $\exists u_0\in X\setminus\{0\}:\ A_\lambda u_0=0$.

\textbf{Доказательство, идея.}
\begin{enumerate}
    \item По билету 20 $\text{Im}\,A_\lambda$ замкнут, поэтому по аннуляторным соотношениям $\text{Im}\,A_\lambda={}^{\perp}\ker(A_\lambda)^*={}^{\perp}\ker A_\lambda^*$.

    \item Поэтому $\exists u\in X:\ A_\lambda u=v$ тогда и только тогда, когда $v\in\text{Im}\,A_\lambda$, что равносильно $v\in{}^{\perp}\ker A_\lambda^*$.

    \item Последнее означает ровно $\forall f\in\ker A_\lambda^*\quad f(v)=0$, что и даёт критерий Фредгольма.

    \item Для альтернативы: если $\ker A_\lambda=0$, то по равенству размерностей ядер получаем $\ker A_\lambda^*=0$.

    \item Тогда $\text{Im}\,A_\lambda={}^{\perp}\ker A_\lambda^*={}^{\perp}0=X$, то есть $\forall v\in X\ \exists u\in X:\ A_\lambda u=v$.

    \item Так как $\ker A_\lambda=0$, это решение единственно: $\forall v\in X\ \exists!u\in X:\ A_\lambda u=v$. Получаем первую ветвь альтернативы.

    \item Если же $\ker A_\lambda\ne0$, то $\exists u_0\in X\setminus\{0\}:\ A_\lambda u_0=0$. Получаем вторую ветвь альтернативы.
\end{enumerate}

\section*{24. Спектр компактного оператора}
\textbf{Формулировка.} Пусть $X$ --- банахово пространство, $A\in\mathcal L(X)$ компактен. Тогда $\forall \lambda\in\sigma(A)\setminus\{0\}$ число $\lambda$ является собственным числом $A$, то есть $\ker A_\lambda\ne0$. Кроме того, $\forall r>0$ множество $\{\lambda\in\sigma(A):|\lambda|>r\}$ конечно, поэтому $\sigma(A)$ не более чем счётен, если счётен, то $\lambda_k \rightarrow 0$, если $X$ бесконечномерно, то $0\in\sigma(A)$.

\textbf{Доказательство, идея.}
\begin{enumerate}
    \item Пусть $\lambda\ne0$ и $\lambda\in\sigma(A)$. Если $\ker A_\lambda=0$, то по альтернативе Фредгольма $\forall y\in X\ \exists!x\in X:\ A_\lambda x=y$, то есть $A_\lambda$ обратим, противоречие с $\lambda\in\sigma(A)$. Значит $\sigma(A) = \sigma_p(A)$.

    \item Докажем, что $\forall r>0$ множество $\{\lambda\in\sigma(A):|\lambda|>r\}$ конечно. Пусть нет: есть собственные числа $\{ \lambda_n \}$, $|\lambda_n|>r$.

    \item Берём собственные векторы $x_n\ne0$, $Ax_n=\lambda_n x_n$, и кладём $L_n=\Lin\{x_1,\ldots,x_n\}$. Собственные векторы, отвечающие различным собственным числам, линейно независимы, значит $L_1\subsetneq L_2\subsetneq\ldots$.

    \item По лемме Рисса выбираем $y_n\in L_n$, $\|y_n\|=1$, так что $\rho(y_n,L_{n-1})>\frac12$.

    \item При $m<n$ имеем $Ay_m\in L_m\subset L_{n-1}$ и $Ay_n-\lambda_n y_n\in L_{n-1}$, поэтому $\|Ay_n-Ay_m\|\ge |\lambda_n|\rho(y_n,L_{n-1})>\frac r2$.

    \item Значит у последовательности $Ay_n$ нет фундаментальной подпоследовательности, что противоречит компактности $A$, так как $\|y_n\|=1$.

    \item Следовательно, $\forall r>0$ множество $\{\lambda\in\sigma(A):|\lambda|>r\}$ конечно. Тогда $\sigma(A)\setminus\{0\}=\bigcup_{n=1}^{\infty}\{\lambda\in\sigma(A):|\lambda|>\frac1n\}$ не более чем счётно, и все возможные точки накопления могут быть только в $0$.

    \item Наконец, $0\in\sigma(A)$: иначе $A^{-1}\in\mathcal L(X)$, и тогда $I=A^{-1}A$ компактен. Но компактность $I$ означает компактность единичного шара, что в бесконечномерном нормированном пространстве невозможно. Противоречие.
\end{enumerate}

\section*{25. Самосопряжённый оператор: вещественность спектра и отсутствие остаточного}
\textbf{Формулировка.} Если $A=A^\ast$ в гильбертовом $\mathcal H$, то $\sigma(A)\subset\R$, а остаточный спектр пуст.

\textbf{Доказательство, идея.}
\begin{enumerate}
    \item Для $\lambda=\mu+i\nu$, $\nu\ne0$, имеем $\|(A-\lambda I)x\|^2=\|(A-\mu I)x\|^2+\nu^2\|x\|^2$, потому что $A-\mu I$ самосопряжён.
    \item Отсюда оценка снизу $\|(A-\lambda I)x\|\ge|\nu|\|x\|$: ядро нулевое, образ замкнут, обратный на образе ограничен.
    \item У сопряжённого $(A-\lambda I)^\ast=A-\bar\lambda I$ тоже нулевое ядро; значит ортогональное дополнение образа нулевое, образ плотен.
    \item Образ одновременно замкнут и плотен, значит весь $\mathcal H$; $\lambda$ в резольвенте.
    \item Для вещественного $\lambda$: если $\Ker(A-\lambda I)=0$, то $(\operatorname{Im}(A-\lambda I))^\perp=\Ker(A-\lambda I)=0$, значит образ плотен. Поэтому остаточного спектра нет.
\end{enumerate}

\section*{26. Критерий спектра самосопряжённого и локализация}
\textbf{Формулировка.} Для $A=A^\ast$ и $\lambda\in\R$: $\lambda\in\sigma(A)$ iff существуют $\|x_n\|=1$ такие, что $\|(A-\lambda I)x_n\|\to0$. Если $m_-=\inf_{\|x\|=1}\langle Ax,x\rangle$, $m_+=\sup_{\|x\|=1}\langle Ax,x\rangle$, то $\sigma(A)\subset[m_-,m_+]$ и $m_\pm\in\sigma(A)$.

\textbf{Доказательство, идея.}
\begin{enumerate}
    \item Для самосопряжённого $A-\lambda I$ вещественная $\lambda$ лежит в резольвенте iff есть оценка снизу $\|(A-\lambda I)x\|\ge c\|x\|$.
    \item Отрицание этой оценки после нормировки даёт последовательность почти собственных векторов.
    \item Если $\lambda>m_+$, то $\langle(\lambda I-A)x,x\rangle\ge(\lambda-m_+)\|x\|^2$, значит по КБШ $\|(A-\lambda I)x\|\ge(\lambda-m_+)\|x\|$; аналогично слева от $m_-$.
    \item Для $m_+$ берём максимизирующую последовательность $x_n$. Псевдоскалярное произведение $\langle u,v\rangle_+=\langle(m_+I-A)u,v\rangle$ и КБШ для него дают $\|(A-m_+I)x_n\|\to0$.
    \item То же для $m_-$.
\end{enumerate}

\section*{27. Теорема Гильберта--Шмидта}
\textbf{Формулировка.} Если $\mathcal H$ бесконечномерно, $A=A^\ast$ компактен, то в $(\Ker A)^\perp$ есть ортонормированный базис из собственных векторов, и $\mathcal H=\Ker A\oplus\bigoplus_n L_n$, где $L_n=\Ker(A-\lambda_n I)$ конечномерны.

\textbf{Доказательство, идея.}
\begin{enumerate}
    \item Спектр компактного самосопряжённого: ненулевые точки --- вещественные собственные числа конечной кратности, возможное накопление только $0$.
    \item В каждом $L_n$ выбираем ОНБ; разные $L_n$ ортогональны: $\lambda_n\langle x,y\rangle=\langle Ax,y\rangle=\langle x,Ay\rangle=\lambda_m\langle x,y\rangle$.
    \item Пусть $L=\overline{\bigoplus_n L_n}$. Тогда $L\subset(\Ker A)^\perp$.
    \item Рассматриваем $M=(\Ker A)^\perp\ominus L$. Оно инвариантно относительно $A$: самосопряжённость переносит ортогональность.
    \item Если $A|_M\ne0$, у него есть ненулевое спектральное значение, значит собственный вектор вне $L$; противоречие. Значит $A|_M=0$, но $M\subset(\Ker A)^\perp$, значит $M=0$.
\end{enumerate}

\section*{28. Резольвентное множество и компактность спектра}
\textbf{Формулировка.} Для $A\in\mathcal L(\X)$ в нетривиальном банаховом $\X$ резольвентное множество открыто, спектр замкнут, ограничен, непуст и потому компактен.

\textbf{Доказательство, идея.}
\begin{enumerate}
    \item Лемма Неймана: если $\|B\|<1$, то $(I-B)^{-1}=\sum_{k\ge0}B^k$.
    \item Если $\lambda\in\rho(A)$, то $A-(\lambda+h)I=(A-\lambda I)(I-h(A-\lambda I)^{-1})$; при малом $h$ второй множитель обратим по Нейману.
    \item Если $|\lambda|>\|A\|$, то $A-\lambda I=-\lambda(I-A/\lambda)$ обратим; значит $\sigma(A)\subset\{|\lambda|\le\|A\|\}$.
    \item Если спектр пуст, резольвента $R_A(\lambda)$ целая; при больших $\lambda$ ряд Неймана даёт $R_A(\lambda)\to0$.
    \item Для $x\in\X$, $f\in\X^\ast$ функция $f(R_A(\lambda)x)$ целая и убывает на бесконечности; по Лиувиллю ноль. Тогда $R_A(\lambda)=0$, невозможно.
\end{enumerate}

\section*{29. Спектральный радиус}
\textbf{Формулировка.} Для $A\in\mathcal L(\X)$ в банаховом пространстве $r(A)=\max_{\lambda\in\sigma(A)}|\lambda|=\lim_{n\to\infty}\|A^n\|^{1/n}$. Кроме того, $r(A)=\|A\|$ iff $\|A^n\|=\|A\|^n$ для всех $n$.

\textbf{Доказательство, идея.}
\begin{enumerate}
    \item Существование предела: $\log\|A^n\|$ субаддитивна; применяем лемму Фекете. Нильпотентный случай отдельно тривиален.
    \item Оценка $r(A)\le L$ следует из спектрального отображения/оценки: $|\lambda|^n\le\|A^n\|$ для $\lambda\in\sigma(A)$, затем корень и предел.
    \item Для $L\le r(A)$ берём $R>r(A)$. На окружности $|\lambda|=R$ резольвента ограничена.
    \item Из ряда Неймана на большой окружности и формулы Коши получаем $f(A^nx)=\frac1{2\pi i}\int_{|\lambda|=R}\lambda^n f(R_A(\lambda)x)\,d\lambda$.
    \item Значит $\|A^n\|\le C_RR^{n+1}$, откуда $L\le R$; пускаем $R\downarrow r(A)$.
    \item Критерий: если $r(A)=\|A\|$, то $\|A\|^n=r(A)^n=r(A^n)\le\|A^n\|\le\|A\|^n$. Обратно сразу из формулы Гельфанда.
\end{enumerate}

\end{document}
